% =========================
% report/implementation.tex
% =========================

\subsection{اصول طراحی پیاده‌سازی}
پیاده‌سازی به‌گونه‌ای طراحی می‌شود که:
\begin{itemize}
  \item \textbf{ماژولار} باشد (هر مرحله مستقل و قابل تست)،
  \item \textbf{بازتولیدپذیر} باشد (ثبت کامل پیکربندی و خروجی‌ها)،
  \item \textbf{قابل مشاهده و قابل تحلیل} باشد (ذخیره خروجی‌های میانی و دیباگ).
\end{itemize}

\subsection{ماژول‌های اصلی (سطح مفهومی)}
\paragraph{(۱) ورودی/خروجی و پیکربندی}
این بخش زوج‌های تصویر را می‌خواند و پارامترها را از فایل تنظیمات دریافت می‌کند. یک پیکربندی علمی باید شامل موارد زیر باشد:
\begin{itemize}
  \item سقف اندازه تصویر $M$،
  \item پارامترهای ORB (مثل \texttt{nfeatures} و تعداد سطوح هرم)،
  \item نوع تطبیق‌گر و قواعد پالایش تطبیق‌ها،
  \item پارامترهای RANSAC مثل آستانه $\tau$،
  \item ضریب هم‌پوشانی $\alpha$،
  \item آستانه‌های ارزیابی (مثل تحمل موفقیت $\delta$).
\end{itemize}

\paragraph{(۲) استخراج ویژگی}
ORB \cite{Rublee2011ORB} نقاط کلیدی و توصیف‌گر دودویی را می‌سازد. انتخاب ORB به دلیل سرعت و کارایی مناسب در کاربردهای مهندسی است.

\paragraph{(۳) تطبیق و پالایش}
توصیف‌گرهای دودویی با فاصلهٔ همینگ تطبیق می‌شوند. سپس با Cross-check و فیلتر فاصله، تطبیق‌های مبهم یا غلط حذف می‌شوند. این کار معمولاً «دقت» را افزایش می‌دهد اما ممکن است «بازخوانی» را کاهش دهد؛ به همین دلیل باید در آزمایش‌های ابلیشن بررسی شود.

\paragraph{(۴) تخمین هموگرافی}
هموگرافی با RANSAC \cite{FischlerBolles1981} تخمین زده می‌شود تا اثر پرت‌ها کاهش یابد. اگر نسبت درون‌نهشتی‌ها پایین باشد، RANSAC باید تکرار بیشتری انجام دهد و هزینه زمانی افزایش می‌یابد؛ بنابراین پالایش تطبیق‌ها هم برای صحت و هم برای سرعت مهم است.

\paragraph{(۵) وارپ و هم‌پوشانی}
الگو با $\mathbf{H}$ به فضای عکس وارپ می‌شود (پرسپکتیو). سپس ترکیب آلفا:
\[
\mathcal{I}_{\text{out}} = \alpha \mathcal{I}_{\text{warp}} + (1-\alpha)\mathcal{I}_Q
\]
خروجی نهایی را می‌سازد.

\subsection{خروجی‌های دیباگ (برای تشخیص علت شکست)}
برای اینکه گزارش علمی «قابل دفاع» باشد، پیاده‌سازی بهتر است موارد زیر را ذخیره کند:
\begin{itemize}
  \item \textbf{تصویر تطبیق خام:} نمایش بهترین $K$ تطبیق قبل از RANSAC.
  \item \textbf{تصویر تطبیق درون‌نهشتی:} نمایش تطبیق‌هایی که با $\mathbf{H}$ سازگارند.
  \item \textbf{تصویر وارپ‌شده:} فقط الگوی وارپ‌شده بدون هم‌پوشانی برای بررسی اعوجاج‌ها.
  \item \textbf{نمایش گوشه‌ها:} فرافکنی چهار گوشهٔ الگو روی عکس (تشخیص «وارپ اشتباه اما منسجم»).
\end{itemize}

\subsection{فایل‌های خروجی و نقش آنها}
\paragraph{هم‌پوشانی‌ها}
\begin{itemize}
  \item \texttt{../results/overlays/*.png}: خروجی اصلی کیفی.
\end{itemize}

\paragraph{معیارها}
\begin{itemize}
  \item \texttt{../results/metrics/per\_image.csv}: معیارهای نمونه‌به‌نمونه (تعداد تطبیق، تعداد درون‌نهشتی، نسبت درون‌نهشتی، خطا، زمان).
  \item \texttt{../results/metrics/summary.csv}: خلاصهٔ سطح داده (میانگین/میانه/صدک‌ها).
  \item \texttt{../results/metrics/config.json}: ثبت دقیق پارامترها برای بازتولید.
\end{itemize}

\paragraph{شکل‌های گزارش}
\begin{itemize}
  \item \texttt{../results/figures/overlay\_grid.png}: شبکه‌ای از هم‌پوشانی‌های نماینده.
  \item \texttt{../results/figures/reproj\_error\_hist.png}: هیستوگرام خطای بازفرافکنی.
  \item \texttt{../results/figures/runtime\_breakdown.png}: تحلیل گلوگاه زمانی.
  \item \texttt{../results/figures/ablation\_plot.png}: حساسیت به پارامترها.
  \item \texttt{../results/figures/failure\_case\_*.png}: نمونه‌های شکست.
\end{itemize}

\subsection{الزامات ثبت (Logging) برای گزارش علمی}
برای نتیجه‌گیری علمی، باید حتماً ثبت شود:
\begin{itemize}
  \item سیاست تغییر اندازه و رزولوشن نهایی،
  \item تعداد نقاط کلیدی در الگو و عکس،
  \item تعداد تطبیق‌ها $N_m$ و تعداد درون‌نهشتی‌ها $N_i$،
  \item نسبت درون‌نهشتی $N_i/N_m$،
  \item (در صورت زمین‌حقیقت) توزیع خطای بازفرافکنی،
  \item زمان اجرا (میانه و صدک ۹۵ توصیه می‌شود).
\end{itemize}

\subsection{پیچیدگی محاسباتی و عملکرد عملی}
اگر $n_T$ و $n_Q$ تعداد نقاط کلیدی باشند، تطبیق خام در بدترین حالت می‌تواند $O(n_T n_Q)$ باشد. در عمل با کنترل \texttt{nfeatures} و اندازه تصویر $M$ زمان اجرا مدیریت می‌شود. هزینه RANSAC نیز به تعداد تکرارها وابسته است و با کاهش نسبت درون‌نهشتی‌ها افزایش می‌یابد \cite{FischlerBolles1981}. بنابراین انتخاب پارامترها باید هم «دقت» و هم «زمان» را همزمان لحاظ کند.

\subsection{درج کُد در گزارش (اختیاری)}
اگر قصد دارید کُد پایتون را داخل PDF نمایش دهید، فایل‌ها را در \texttt{../src/} قرار دهید و از دستور زیر استفاده کنید:
\begin{verbatim}
\SafeListing{../src/alignment.py}{ماژول هم‌ترازی (ویژگی‌ها، تطبیق، هموگرافی).}{lst:alignment}
\end{verbatim}
